\section{Phân tích log file [Nâng cao]}

\subsection{Phân tích \& Thiết kế (M0)}
\begin{itemize}
    \item \textbf{Input:}
    \begin{itemize}
        \item File \texttt{log.txt}, mỗi dòng có dạng:
        \begin{center}
        \texttt{YYYY-MM-DD HH:MM:SS [LEVEL] message}
        \end{center}
        với LEVEL $\in$ \{\texttt{INFO, WARN, ERROR}\}
    \end{itemize}

    \item \textbf{Xử lý:}
    \begin{itemize}
        \item \textbf{Case biên:}
        \begin{itemize}
            \item File rỗng $\rightarrow$ tất cả count = 0, \texttt{first\_error = NONE}, \texttt{last\_error = NONE}
            \item Dòng sai format $\rightarrow$ bỏ qua
            \item Không có dòng ERROR $\rightarrow$ \texttt{first\_error = NONE}, \texttt{last\_error = NONE}
        \end{itemize}

        \item \textbf{Module 1 (Đọc file):}
        \begin{itemize}
            \item Hàm \texttt{readLines(path, vector<string>\&)}
            \item Đọc toàn bộ dòng bằng \texttt{getline}
        \end{itemize}

        \item \textbf{Module 2 (Parse dòng log):}
        \begin{itemize}
            \item Hàm \texttt{parseLog(line, timestamp, level)}
            \item Tách:
            \begin{itemize}
                \item \texttt{timestamp}: 19 ký tự đầu
                \item \texttt{level}: nằm trong dấu \texttt{[ ]}
            \end{itemize}
            \item Nếu không đúng format $\rightarrow$ return false
        \end{itemize}

        \item \textbf{Module 3 (Thống kê):}
        \begin{itemize}
            \item Đếm số dòng theo từng LEVEL:
            \[
            \texttt{INFO}, \texttt{WARN}, \texttt{ERROR}
            \]
            \item Nếu gặp ERROR:
            \begin{itemize}
                \item Nếu là ERROR đầu tiên $\rightarrow$ lưu \texttt{first\_error}
                \item Luôn cập nhật \texttt{last\_error}
            \end{itemize}
        \end{itemize}

        \item \textbf{Module 4 (Ghi file):}
        \begin{itemize}
            \item Hàm \texttt{writeSummary(path, info, warn, error, first, last)}
        \end{itemize}
    \end{itemize}

    \item \textbf{Output:}
    \begin{center}
    \texttt{INFO=<n> / WARN=<n> / ERROR=<n> / first\_error=<timestamp> / last\_error=<timestamp>}
    \end{center}
\end{itemize}

\subsection{Mã nguồn C++11 (Tách module)}
\begin{lstlisting}
#include <iostream>
#include <fstream>
#include <vector>
#include <string>

using namespace std;

// MODULE 1: Read file
bool readLines(string path, vector<string>& lines) {
    ifstream file(path);
    if (!file.is_open()) return false;

    string line;
    while (getline(file, line)) {
        lines.push_back(line);
    }

    file.close();
    return true;
}

// MODULE 2: Parse log line
bool parseLog(const string& line, string& timestamp, string& level) {
    if (line.size() < 22) return false;

    // timestamp = 19 ky tu dau
    timestamp = line.substr(0, 19);

    // tim [LEVEL]
    size_t l = line.find('[');
    size_t r = line.find(']');

    if (l == string::npos || r == string::npos || r <= l + 1) return false;

    level = line.substr(l + 1, r - l - 1);

    if (level != "INFO" && level != "WARN" && level != "ERROR")
        return false;

    return true;
}

// MODULE 3: Analyze
void analyzeLogs(const vector<string>& lines,
                 int& info, int& warn, int& error,
                 string& first_error, string& last_error) {

    info = warn = error = 0;
    first_error = "NONE";
    last_error = "NONE";

    for (const string& line : lines) {
        string timestamp, level;

        if (!parseLog(line, timestamp, level)) continue;

        if (level == "INFO") info++;
        else if (level == "WARN") warn++;
        else if (level == "ERROR") {
            error++;
            if (first_error == "NONE") first_error = timestamp;
            last_error = timestamp;
        }
    }
}

// MODULE 4: Write output
bool writeSummary(string path,
                  int info, int warn, int error,
                  const string& first_error,
                  const string& last_error) {

    ofstream file(path);
    if (!file.is_open()) return false;

    file << "INFO=" << info
         << " / WARN=" << warn
         << " / ERROR=" << error
         << " / first_error=" << first_error
         << " / last_error=" << last_error;

    file.close();
    return true;
}

// MAIN
int main() {
    vector<string> lines;

    if (!readLines("log.txt", lines)) {
        cerr << "Khong mo duoc file log.txt\n";
        return 1;
    }

    int info, warn, error;
    string first_error, last_error;

    analyzeLogs(lines, info, warn, error, first_error, last_error);

    if (!writeSummary("summary.txt", info, warn, error, first_error, last_error)) {
        cerr << "Khong ghi duoc file summary.txt\n";
        return 1;
    }

    return 0;
}
\end{lstlisting}

\subsection{Kế hoạch kiểm thử (Test Cases)}

\textbf{Test Case 1 (Thông thường)}

\textbf{Input (log.txt):}
\begin{lstlisting}
2024-01-01 10:00:00 [INFO] Start system
2024-01-01 10:01:00 [WARN] Low memory
2024-01-01 10:02:00 [ERROR] Crash detected
2024-01-01 10:03:00 [ERROR] Restart failed
\end{lstlisting}

\textbf{Expected Output:}
\begin{lstlisting}
INFO=1 / WARN=1 / ERROR=2 / first_error=2024-01-01 10:02:00 / last_error=2024-01-01 10:03:00
\end{lstlisting}

\vspace{0.5cm}

\textbf{Test Case 2 (Không có ERROR)}

\textbf{Input (log.txt):}
\begin{lstlisting}
2024-01-01 10:00:00 [INFO] Start
2024-01-01 10:01:00 [WARN] Something
\end{lstlisting}

\textbf{Expected Output:}
\begin{lstlisting}
INFO=1 / WARN=1 / ERROR=0 / first_error=NONE / last_error=NONE
\end{lstlisting}

\vspace{0.5cm}

\textbf{Test Case 3 (Case biên -- file rỗng)}

\textbf{Input (log.txt):}
\begin{lstlisting}
\end{lstlisting}

\textbf{Expected Output:}
\begin{lstlisting}
INFO=0 / WARN=0 / ERROR=0 / first_error=NONE / last_error=NONE
\end{lstlisting}